\begin{prob}
    Roughly how long will it take for a linear time algorithm to run? What
    about a quadratic time algorithm? Or worse, a cubic? In this problem,
    we'll estimate these times.

    % Write the problem here.
    Suppose algorithm A takes $n$ microseconds to run on a problem of size $n$,
    while algorithm B takes $n^2$ microseconds and algorithm C takes $n^3$
    microseconds (recall that a microsecond is one millionth of a second).
    How long will each algorithm take to run when the input is of size one
    thousand, ten thousand, one hundred thousand, and one million? That is,
    fill in the following table:

    \begin{center}
    \begin{tabular}{rcccc}
        \toprule
        & $n = 1{,}000$ & $n = 10{,}000$ & $n = 100{,}000$ & $n = 1{,}000{,}000$ \\ \midrule
        A (Linear) & 0.00 s & 0.01 s & 0.10 s & 1 s \\
        B (Quadratic) & ? & ? & ? & ? \\
        C (Cubic) & ? & ? & ? & ? \\\bottomrule
    \end{tabular}
    \end{center}

    The answers for Algorithm A are already provided; you can use them to check your strategy.

    Express each time in either seconds, minutes, hours, days, or years. Use
    the largest unit that you can without getting an answer less than one. For
    example, instead of ``365 days'', say ``1 year''; but use ``364 days''
    instead of ``0.997 years''. Round to two decimal places (it's OK for an
    answer to round to 0.00).

    Hint: you can calculate your answers by hand, or you can write some code to compute them.
    If you write code, provide it with your solution -- if you solve by hand, show your calculations.

    \begin{soln}

        \begin{center}
        \begin{tabular}{rrrrr}
            \toprule
            & $n = 1{,}000$ & $n = 10{,}000$ & $n = 100{,}000$ & $n = 1{,}000{,}000$ \\ \midrule
            A & 0.00 s & 0.01 s & 0.10 s & 1 s \\
            B & 1 s & 1.67 m & 2.78 h & 11.57 d \\
            C & 16.67 m & 11.57 d & 31.71 y & 31{,}709.79 y \\\bottomrule
        \end{tabular}
        \end{center}

        This was computed by the following code:

        \inputminted{python}{\thisdir/include-solution/code.py}

    \end{soln}

\end{prob}
